% appendices/appendix-a-rubrik.tex
\chapter{Rubrik Penilaian Lengkap}
\label{app:rubrik}

Lampiran ini memuat rubrik penilaian dalam bentuk lengkap, yaitu beserta
uraian setiap tingkat capaian. Bab~XVI dan Bab~XVII telah menyajikan bobot dan
kriteria pokoknya; yang ditambahkan di sini adalah perbedaan antara tingkat
yang satu dan tingkat berikutnya, agar penilaian tidak bergantung pada kesan
penilai.

\section{Cara Membaca Rubrik}

Setiap kriteria dinilai pada empat tingkat. Tingkat yang dipakai adalah
tingkat terendah yang \emph{seluruh} syaratnya terpenuhi, bukan tingkat
tertinggi yang sebagian syaratnya terpenuhi. Apabila sebuah pekerjaan memenuhi
tiga dari empat syarat pada tingkat ``Sangat Baik'', ia dinilai ``Baik''.

\begin{table}[htbp]
\centering
\caption{Empat tingkat capaian yang dipakai pada seluruh rubrik}
\label{tab:tingkat-capaian}
{\setlength{\tabcolsep}{4pt}\footnotesize%
\begin{tabularx}{\textwidth}{@{}L{3.4cm} C{2.2cm} Y@{}}
\toprule
\headrow \textbf{Tingkat} & \textbf{Rentang nilai} & \textbf{Ciri umum} \\
\midrule
Sangat Baik & 85--100 & Seluruh syarat terpenuhi, dan pekerjaannya dapat dipakai
sebagai contoh bagi kelompok lain. \\
\addlinespace
Baik & 70--84 & Seluruh syarat pokok terpenuhi, dengan satu kekurangan kecil
yang tidak mengubah hasil. \\
\addlinespace
Cukup & 60--69 & Syarat pokok terpenuhi sebagian; hasilnya berjalan, tetapi
masih menuntut perbaikan sebelum dapat dipakai. \\
\addlinespace
Kurang & $<$ 60 & Syarat pokok tidak terpenuhi; hasilnya tidak berjalan atau
tidak dapat diperiksa. \\
\bottomrule
\end{tabularx}}
\end{table}

\section{Rubrik Mini Proyek AI}

Mini Proyek AI menilai \CPMK{2}, yaitu perilaku NPC. Bobotnya 10 persen dari
nilai akhir, dan dikerjakan berkelompok.

\begin{landscape}
\begin{table}[htbp]
\centering
\caption{Rubrik Mini Proyek AI beserta uraian empat tingkat capaian}
\label{tab:rubrik-mini-lengkap}
{\setlength{\tabcolsep}{4pt}\footnotesize%
\begin{tabularx}{\linewidth}{@{}L{2.8cm} C{1.3cm} Y Y Y Y@{}}
\toprule
\headrow \textbf{Kriteria} & \textbf{Bobot} & \textbf{Sangat Baik} & \textbf{Baik} & \textbf{Cukup} & \textbf{Kurang} \\
\midrule
Implementasi FSM & 25\% &
Keadaan terpisah dan jelas; setiap transisi berjalan menurut jarak dan
\istilah{cooldown}; tidak ada keadaan yang tidak dapat ditinggalkan. &
Transisi berjalan dengan benar; satu keadaan kurang lengkap, misalnya tidak
memiliki jalan kembali. &
Sebagian transisi tidak berjalan; NPC sempat diam tanpa alasan yang jelas. &
Keadaan tidak terpisah; perilaku ditulis sebagai rangkaian \code{if} yang
saling bertumpuk. \\
\addlinespace
Implementasi A* & 25\% &
Jalur benar pada peta berhalangan; memakai antrean prioritas; heuristik
dipilih dengan alasan yang dapat dipertanggungjawabkan. &
Jalur benar dan memakai antrean prioritas; heuristik dipakai tanpa penjelasan. &
Jalur benar pada peta sederhana, tetapi gagal atau lambat pada peta
berhalangan. &
Jalur tidak selalu benar, atau pencarian tidak berakhir pada peta tertentu. \\
\addlinespace
Integrasi FSM dan A* & 20\% &
Pencarian dipicu tepat pada keadaan pengejar; pencarian ulang berkala; kinerja
terjaga dan alasannya dijelaskan. &
Pencarian dipicu pada keadaan yang tepat, tanpa pencarian ulang berkala. &
Pencarian dipicu pada keadaan yang kurang tepat, sehingga jalurnya tertinggal
dari pergerakan pemain. &
Pencarian dijalankan setiap bingkai, sehingga permainan melambat. \\
\addlinespace
Pengujian dan laporan & 20\% &
Sekurang-kurangnya tiga skenario uji; pengamatan objektif dengan angka;
perbaikan ditindaklanjuti dan hasilnya dicatat. &
Tiga skenario uji dan pengamatan objektif; perbaikan belum ditindaklanjuti. &
Skenario uji kurang dari tiga, atau pengamatannya berupa kesan. &
Tidak ada skenario uji maupun laporan. \\
\addlinespace
Presentasi dan kerja sama & 10\% &
Demonstrasi jelas dan tepat waktu; setiap anggota memahami bagiannya dan dapat
menjawab pertanyaan. &
Demonstrasi jelas; satu anggota kurang menguasai bagiannya. &
Demonstrasi bergantung pada satu orang. &
Demonstrasi tidak disiapkan. \\
\bottomrule
\end{tabularx}}
\end{table}
\end{landscape}

\section{Rubrik Proyek Akhir}

Proyek Akhir menilai \CPMK{3} dan \CPMK{4}, dengan bobot 30 persen dari nilai
akhir. Rubrik ini adalah instrumen terbesar pada mata kuliah ini, dan
satu-satunya yang menilai kemampuan menggabungkan bagian-bagian.

\begin{landscape}
\begin{table}[htbp]
\centering
\caption{Rubrik Proyek Akhir beserta uraian empat tingkat capaian}
\label{tab:rubrik-akhir-lengkap}
{\setlength{\tabcolsep}{4pt}\footnotesize%
\begin{tabularx}{\linewidth}{@{}L{2.8cm} C{1.3cm} Y Y Y Y@{}}
\toprule
\headrow \textbf{Kriteria} & \textbf{Bobot} & \textbf{Sangat Baik} & \textbf{Baik} & \textbf{Cukup} & \textbf{Kurang} \\
\midrule
Fungsionalitas mekanik inti dan level & 30\% &
Level selesai dari awal sampai akhir tanpa penghambat; mekanik inti bekerja
sesuai rancangan; kendali pemain responsif; tidak ditemukan
\istilah{soft-lock}. &
Level selesai, dan mekanik intinya bekerja; satu mekanik tambahan belum
berjalan. &
Level dapat dimainkan, tetapi ada bagian yang membingungkan atau satu tempat
yang hanya dapat dicapai dengan susah payah. &
Level tidak dapat diselesaikan, atau mekanik intinya tidak bekerja. \\
\addlinespace
Integrasi aset visual, animasi, dan UI & 25\% &
Aset terpasang rapi dan seragam ukurannya; animasi berganti menurut keadaan;
antarmuka menampilkan angka yang selalu benar. &
Ketiganya terpasang; satu keadaan animasi atau satu angka antarmuka belum
tepat. &
Aset terpasang, tetapi animasi atau antarmuka belum terhubung sepenuhnya pada
keadaan permainan. &
Aset belum terpasang, atau antarmuka menampilkan angka yang tidak sesuai
keadaan. \\
\addlinespace
Integrasi audio dan umpan balik & 15\% &
Musik latar dan efek suara terpisah; setiap efek suara terpicu pada
peristiwanya; tingkat kekuatan beralasan; ada pasangan yang terlihat. &
Efek suara terpicu dengan benar; tingkat kekuatan belum ditata, atau satu
peristiwa belum memperoleh bunyi. &
Sebagian bunyi terpicu, atau musik tertutup oleh efek suara. &
Bunyi belum terpasang, atau terdengar pada peristiwa yang salah. \\
\addlinespace
Kualitas kode dan struktur \istilah{project} & 15\% &
\istilah{Project} tersusun rapi; penamaan konsisten; tidak ada kode mati;
bagian yang berulang disatukan. &
Struktur rapi dan penamaan konsisten; ada sedikit kode mati. &
Struktur dapat dipahami, tetapi penamaannya tidak konsisten atau ada kode yang
tidak terpakai. &
Seluruh kode berada pada satu berkas, atau strukturnya tidak dapat dipahami. \\
\addlinespace
\istilah{Playtest}, laporan, dan presentasi & 15\% &
Lembar \istilah{playtest} terisi oleh dua penguji atau lebih; \istilah{log bug}
lengkap dengan status; laporan memuat keadaan sebelum dan sesudah; demo tepat
waktu. &
Kelengkapan lengkap; satu bagian laporan belum memuat bukti. &
Kelengkapan sebagian; \istilah{log bug} tidak memuat status penanganan. &
Kelengkapan tidak diserahkan, atau demo tidak dilakukan. \\
\bottomrule
\end{tabularx}}
\end{table}
\end{landscape}

\section{Rubrik Portofolio Praktikum}

Tugas praktikum dan portofolio menilai \CPMK{1} sampai \CPMK{4} dengan bobot
25 persen. Penilaiannya dilakukan atas artefak mingguan pada repositori,
bukan atas satu pekerjaan besar.

\begin{table}[htbp]
\centering
\caption{Rubrik portofolio praktikum}
\label{tab:rubrik-portofolio}
{\setlength{\tabcolsep}{4pt}\footnotesize%
\begin{tabularx}{\textwidth}{@{}L{3.4cm} C{1.4cm} Y@{}}
\toprule
\headrow \textbf{Kriteria} & \textbf{Bobot} & \textbf{Deskripsi capaian} \\
\midrule
Kelengkapan artefak & 40\% & Seluruh praktikum terkumpul, masing-masing dengan
kode, keluaran, dan catatan singkat. \\
\addlinespace
Kebenaran hasil & 30\% & Kode berjalan dan menghasilkan keluaran yang sesuai
dengan tujuan praktikum. \\
\addlinespace
Catatan dan penjelasan & 20\% & Setiap artefak disertai penjelasan mengapa
caranya dipilih, bukan hanya apa yang dikerjakan. \\
\addlinespace
Riwayat penyimpanan & 10\% & Pekerjaan disimpan bertahap dengan pesan
penyimpanan yang menerangkan isinya. \\
\bottomrule
\end{tabularx}}
\end{table}

\section{Rubrik Laporan dan Refleksi}

Rubrik ini dipakai bersama Proyek Akhir, dan disajikan lengkap pada
Tabel~\ref{tab:rubrik-laporan-lengkap}.

\begin{table}[htbp]
\centering
\caption{Rubrik laporan pengembangan dan refleksi beserta tingkat capaian}
\label{tab:rubrik-laporan-lengkap}
{\setlength{\tabcolsep}{4pt}\footnotesize%
\begin{tabularx}{\textwidth}{@{}L{3.0cm} C{1.2cm} Y Y@{}}
\toprule
\headrow \textbf{Kriteria} & \textbf{Bobot} & \textbf{Sangat Baik atau Baik} & \textbf{Cukup atau Kurang} \\
\midrule
Kelengkapan dokumen & 20\% &
Memuat ringkasan desain, daftar aset beserta lisensinya, lembar
\istilah{playtest}, dan \istilah{log bug} dengan status penanganan. &
Salah satu kelengkapan tidak diserahkan, atau daftar aset tidak memuat
lisensinya. \\
\addlinespace
Laporan perbaikan & 30\% &
Setiap perbaikan disertai keadaan sebelum, keadaan sesudah, dan buktinya. &
Perbaikan diceritakan tanpa bukti, atau hanya sebagian perbaikan yang
dilaporkan. \\
\addlinespace
Kesesuaian GDD dan \istilah{build} & 20\% &
Perbedaan dicantumkan dengan jujur menurut tiga keadaan pada
Tabel~\ref{tab:gdd-build}, beserta alasannya. &
Perbedaan tidak dicantumkan, atau janji dihapus dari GDD agar tampak sesuai. \\
\addlinespace
Refleksi teknologi & 30\% &
Setiap pernyataan memuat keadaan konkret, tindakan, dan hasil; ketiga tinjauan
dijawab terpisah. &
Refleksi berupa kesan umum, atau hanya sebagian tinjauan yang dijawab. \\
\bottomrule
\end{tabularx}}
\end{table}

\section{Ketentuan Kelulusan}

Kelulusan mata kuliah ini ditentukan oleh tiga ketentuan yang berlaku
sekaligus: nilai akhir sekurang-kurangnya 60, kehadiran sekurang-kurangnya
75 persen, dan penyerahan seluruh kelengkapan Proyek Akhir. Bobot setiap
komponen, matriks sumbangannya kepada CPMK, dan konversi nilainya disajikan
pada Bagian~\ref{sec:bobot-nilai}.
